\subsection{COMPONENTS OF A TOPOLOGICAL GROUP}

Let V be a symmetric neighbourhood of e in G. The subgroup generated by V, which is denoted by \(V^{\infty}\) consists of all products \(\prod_{i=1}^{n} x_{i}\) of finite sequences of elements of V. \(V^{\infty}\) is open, since it has e as an interior point, and is therefore closed by Proposition 4 of no. 1. It follows that:

PROPOSITION 6. A connected topological group is generated by every neighbourhood of the identity element.

The converse of this proposition is in general false, as we shall see in Chapter IV (§ 2, no. 5). If a topological group G is generated by each neighbourhood of the identity element, the most we can say is that G contains no open subgroup other than G.

* As an example of a non-connected group G which has an open subgroup distinct from G, we may cite the multiplicative group \(\mathbf{R}^{*}\) of non-zero real numbers, in which the subgroup \(\mathbf{R}_{+}^{*}\) of real numbers \textgreater0 is both open and closed (see Chapter IV, § 3, no. 2). *

PROPOSITION 7. In a topological group G, the component K of the identity element e is a closed normal subgroup. The component of any point \(x \in G\) is the coset \(x.K = K.x\) .

If \(a \in \mathbf{K}\) , then \(a^{-1}\mathbf{K}\) is connected and contains \(e\) ; hence \(\mathbf{K}^{-1}\mathbf{K} \subset \mathbf{K}\) , which shows that \(\mathbf{K}\) is a subgroup of \(\mathbf{G}\) . This subgroup is invariant under all automorphisms of \(\mathbf{G}\) , and in particular under all inner automorphisms, therefore \(\mathbf{K}\) is normal in \(\mathbf{G}\) ; also \(\mathbf{K}\) is closed (Chapter I, § 11, no. 5, Proposition 9). Finally, the left translation \(y \to xy\) is a homeomorphism of \(\mathbf{G}\) which sends \(e\) to \(x\) , and hence the component of \(x\) is \(x\cdot \mathbf{K}\) .

The component of the identity element e of G is called the identity component of G.

